\section{Thermal Processing and Microstructure}
\label{sec:densification}

The thermal step converts a chemically mixed powder into a connected garnet skeleton. It also controls lithium loss, grain growth, residual porosity and the composition of grain boundaries. Because these outcomes are coupled, a sintering temperature alone is not a transferable process descriptor \cite{janani2014influence,moy2025irreversible,chen2022processing}.

\subsection{Experimental conclusion bundle}

The reported experiments support five bounded conclusions. First, sintering additives or lithium-containing transient phases can increase densification and lower resistance, but the additive may remain at boundaries or change the phase assemblage \cite{janani2014influence,luo2019influence,park2026li2co3}. Second, lithium and other elements can diffuse irreversibly during firing, so the compact composition after sintering is not necessarily the weighed composition \cite{moy2025irreversible,gullbrekken2025phase}. Third, rapid microwave or laser schedules can shorten thermal exposure and produce dense LLZO or LLZTO, but they create larger thermal gradients and surface-chemistry contrasts than a conventional furnace \cite{aman2026susceptor,ramos2026ultrafast,zhang2024unveiling}. Fourth, crucible material and cover configuration are chemically active boundary conditions; Pt or Al crucibles should be recorded rather than treated as interchangeable hardware \cite{k2026influence}. Fifth, grain-boundary segregation is not an immutable property of the garnet lattice: co-doping, atmosphere and thermal path can remove or create boundary-rich regions \cite{ahmad2025reducing,golmohammad2024effects,yao2026elimination}.

Recent one-step Al--N doping, temperature-window and transformation-toughening studies broaden the thermal evidence while retaining the same caveat: density, phase assemblage and surface chemistry must be measured together \cite{zhang2026al,qin2026influence,lessmeier2026transformation}. Earlier reactive-sintering and lithium-oxide-additive work reaches high density through transient liquids or enhanced particle contact, providing useful process precedents without proving a universal additive mechanism for every LLZO composition \cite{li2020solid,li2013influence,jonson2018tape,zhang2022high}.

\begin{table}[t]
\centering
\caption{Thermal and forming variables linked to measured structural outcomes.}
\label{tab:thermal}
\begin{tabular}{P{0.20\linewidth}P{0.25\linewidth}P{0.25\linewidth}P{0.20\linewidth}}
\toprule
Process variable & Typical structural consequence & Conductivity interpretation & Representative sources \\
\midrule
Sintering aid or lithium-rich transient & Faster neck growth; possible boundary phase or altered lithium inventory & Total resistance can fall without a bulk change & \cite{janani2014influence,luo2019influence,park2026li2co3} \\
Cover powder, crucible and atmosphere & Changes lithium chemical potential and surface reaction & Phase purity and surface impedance must be checked together & \cite{k2026influence,saran2024optical,saran2025structural} \\
Rapid microwave or laser heating & High density at short dwell; thermal gradients and surface layers & Compare matched density and depth-resolved chemistry & \cite{aman2026susceptor,ramos2026ultrafast,zhang2024unveiling} \\
Co-doping and boundary control & Changes grain size, segregation and boundary continuity & Separate bulk and boundary resistance; do not infer mechanism from total conductivity & \cite{ahmad2025reducing,golmohammad2024effects,yao2026elimination} \\
\bottomrule
\end{tabular}
\end{table}

\subsection{Density, grains and boundaries}

Relative density is the first microstructural prerequisite for any conductivity comparison, but it is not sufficient. Grain size, pore connectivity, boundary thickness and dopant segregation determine the fraction of current that traverses resistive interfaces. Studies combining atomistic and mesoscale descriptions, together with direct boundary imaging, show that a nominally dense pellet can retain heterogeneous transport pathways \cite{heo2021microstructural,gao2021revealing,yao2026elimination}. The minimum structural package is Archimedes density, cross-sectional microscopy, grain-size statistics and a statement on open porosity. Air-stability and microstructure-engineering studies show that this package should include surface carbonate and a depth-resolved chemistry when powders or pellets have seen humid air \cite{nasir2023li,hoinkis2023particle,uhlenbruck2018reactions}.

The same density--microstructure coupling appears in hot pressing, spark-plasma and cold-sintering literature. Grain-boundary state, local pore evolution and the applied current or pressure can change shrinkage and mechanical constraint independently of nominal composition \cite{ni2012room,srivastav2020localized,srivastav2020microstructure,orru2009consolidation,galotta2021the,cai2023microstructure}. These are processing analogues, not LLZO conductivity values, and they motivate reporting the green-body density and the spatial distribution of residual porosity.

\subsection{Membranes and non-bulk form factors}

Tape-cast sheets, bilayer dense--porous membranes and ultrafast-sintered self-standing LLZO demonstrate that the process--structure relation changes when thickness falls below a conventional pellet scale \cite{gao2019preparation,ye2020water,zhang2023bilayer,zhang2023ultrafast,touidjine2026tapex}. Drying stress, binder burnout and in-plane shrinkage can produce gradients that are invisible in a single polished cross-section. Membrane reports should therefore include thickness maps, areal density, bending or handling history and the electrode area used for impedance.

Thin-film and scaffold studies supply complementary form factors: lamellar garnet cells, pulsed-laser-deposited Ta garnets, 3D lithium hosts and porous mixed-conductor interlayers all show that current constriction and mechanical support depend on architecture as well as on the ceramic phase \cite{du2015all,reinacher2014preparation,liu20183d,li20223d,kim2021porous}. The comparison unit is consequently a characterized membrane or composite cross-section, not a powder composition in isolation.

The process map in Fig.~\ref{fig:process-map} is intentionally read from left to right: a powder route does not “cause” a conductivity value directly; it changes phase, porosity and interfaces, which are then sampled by a measurement geometry. The coupled-variable feedbacks are consolidated in Fig.~\ref{fig:coupled-variables} and discussed in Section~\ref{sec:discussion}.
